\documentclass[11pt]{article}
\usepackage{varnernotes}

\begin{document}

\lecturenotes{15a}{Market Crashes and Agent-Based Wolfram Markets}{Fall 2026}{December 1, 2026}{Jeffrey D. Varner}

\learningobjectives
  {Measure abrupt market stress}
  {Distinguish a large return from a drawdown and connect leverage, liquidity, crowded strategies, and feedback to crash amplification.}
  {Encode interacting agents}
  {Represent buy--hold--sell policies as totalistic base-3 rules and update a heterogeneous trader network synchronously.}
  {Evaluate an agent-based model}
  {Map actions into price impact and test whether simulated returns reproduce targeted facts without confusing emergence with causal proof.}

\section{A Crash Is a Path, Not a Single Label}

Markets have no universal numerical definition of a crash. A large one-period loss measures speed; a drawdown measures cumulative loss from the previous peak. For a positive price or wealth index $P_t$, define the log return and running drawdown by
\begin{equation}
  \label{eq:return-drawdown}
  r_t:=\log\!\left(\frac{P_t}{P_{t-1}}\right),
  \qquad
  D_t:=1-\frac{P_t}{M_t},
  \qquad
  M_t:=\max_{0\leq s\leq t}P_s.
\end{equation}
A one-day collapse can make both $r_t$ and $D_t$ extreme; a long bear market can generate a deep $D_t$ without one exceptional daily return. Any empirical ``crash'' threshold must therefore state the index, sampling interval, return convention, peak rule, and recovery rule.

Historical episodes differ. The 1987 U.S. equity decline combined portfolio-insurance selling, market fragmentation, and reduced liquidity. The 2010 Flash Crash unfolded within minutes amid large order flow and interacting automated liquidity, then substantially reversed. The global financial crisis involved leverage, funding stress, opaque exposures, and institutional contagion over months. A shared outcome does not imply one universal cause.

\begin{figure}[ht]
  \centering
  \begin{tikzpicture}[>=Latex,font=\sffamily\small,node distance=1.25cm]
    \node[draw=vnink,very thick,rounded corners,fill=vnsoft,minimum width=2.35cm,minimum height=0.9cm,align=center] (shock) {shock or belief\\revision};
    \node[draw=vnink,very thick,rounded corners,fill=white,minimum width=2.35cm,minimum height=0.9cm,align=center,right=of shock] (sell) {selling and\\margin calls};
    \node[draw=vnink,very thick,rounded corners,fill=vnsoft,minimum width=2.35cm,minimum height=0.9cm,align=center,right=of sell] (liq) {liquidity withdraws\\price impact rises};
    \node[draw=vnink,very thick,rounded corners,fill=white,minimum width=2.35cm,minimum height=0.9cm,align=center,right=of liq] (loss) {lower prices and\\larger losses};
    \draw[vnlink,very thick,->] (shock)--(sell);
    \draw[vnlink,very thick,->] (sell)--(liq);
    \draw[vnlink,very thick,->] (liq)--(loss);
    \draw[vncarnelian,very thick,->,bend left=24] (loss) to node[below,align=center] {leverage and common rules\\amplify the next sale} (sell);
  \end{tikzpicture}
  \caption{A generic crash-amplification loop. The initiating event and market institutions differ across episodes, but endogenous feedback can make selling beget further selling.}
  \label{fig:crash-loop}
\end{figure}

Several mechanisms can reinforce one another:
\begin{itemize}[leftmargin=1.3em,itemsep=2pt,topsep=2pt]
  \item \textbf{Leverage and constraints:} losses reduce equity, trigger margin calls, or force risk-targeting strategies to deleverage.
  \item \textbf{Liquidity mismatch:} desired sales exceed immediate risk-bearing capacity, so price impact and bid--ask spreads increase.
  \item \textbf{Crowding and common information:} similar signals, benchmarks, or stop rules synchronize trades that appeared diversified in normal periods.
  \item \textbf{Network contagion:} exposures, funding links, and imitation transmit distress between agents and institutions.
  \item \textbf{Reflexivity:} falling price changes measured volatility, collateral, beliefs, and allowable positions, feeding back into order flow.
\end{itemize}

Standard diffusion models with continuous paths make sufficiently abrupt moves extremely rare unless volatility is already enormous. Jumps, stochastic volatility, liquidity state variables, and agent interactions offer different ways to represent observed discontinuity and clustering. These are complementary modeling choices, not automatic causal explanations.

\section{Why Build an Agent-Based Market?}

An agent-based model starts from heterogeneous decision rules and an interaction structure, then studies the aggregate dynamics generated by repeated local actions. Unlike a representative-agent equilibrium, it can encode imitation, bounded information, discrete positions, institutional constraints, and explicit network topology. Its strength is mechanism exploration: ``Could these rules generate this pattern?'' Its weakness is non-uniqueness: many micro-models can generate similar macro statistics.

In a Wolfram Market, traders occupy vertices of a grid or graph. Each agent state is also its proposed action,
\begin{equation}
  \label{eq:states}
  s_i(t)\in\{0,1,2\}
  \equiv\{\text{buy},\text{hold},\text{sell}\}.
\end{equation}
Let $\mathcal N_i$ be the $m_i$ neighbors observed by agent $i$. A totalistic rule depends on the neighborhood sum rather than the ordered configuration:
\begin{equation}
  \label{eq:totalistic-state}
  q_i(t):=\sum_{j\in\mathcal N_i}s_j(t)
  \in\{0,1,\ldots,2m_i\},
  \qquad
  s_i(t+1)=\phi_i(q_i(t)).
\end{equation}
Thus a fixed-degree agent needs $2m_i+1$ rule entries, not $3^{m_i}$ entries. The full ordered neighborhood has $3^{m_i}$ possible configurations, but the totalistic map intentionally collapses configurations with the same sum.

For $m$ fixed, write the rule table as digits $d_q:=\phi(q)\in\{0,1,2\}$ and encode it by the base-3 integer
\begin{equation}
  \label{eq:rule-index}
  I_\phi:=\sum_{q=0}^{2m}d_q3^q.
\end{equation}
The digit ordering convention must be documented; reversing it produces a different policy from the same printed digit string.

\begin{figure}[ht]
  \centering
  \begin{tikzpicture}[x=0.72cm,y=0.72cm,font=\sffamily\small]
    \foreach \x in {0,...,6}{
      \foreach \y in {0,...,4}{
        \pgfmathtruncatemacro{\c}{mod(2*\x+\y,3)}
        \ifnum\c=0\def\fc{vnlink!65}\fi
        \ifnum\c=1\def\fc{vnsoft}\fi
        \ifnum\c=2\def\fc{vncarnelian!65}\fi
        \draw[fill=\fc,draw=white] (\x,\y) rectangle ++(1,1);
      }
    }
    \draw[vnink,very thick] (2,1) rectangle ++(3,3);
    \node[anchor=west] at (7.45,3.75) {local neighborhood};
    \node[anchor=west,text=vnmuted] at (7.45,2.85) {$q_i(t)=\sum_{j\in\mathcal N_i}s_j(t)$};
    \node[anchor=west,text=vnlink] at (7.45,1.95) {$s_i(t+1)=\phi_i(q_i(t))$};
    \node[anchor=west] at (7.45,0.8) {buy \quad hold \quad sell};
    \fill[vnlink!65] (7.45,0.15) rectangle ++(0.55,0.45);
    \fill[vnsoft] (8.75,0.15) rectangle ++(0.55,0.45);
    \fill[vncarnelian!65] (10.0,0.15) rectangle ++(0.55,0.45);
  \end{tikzpicture}
  \caption{A three-state cellular market. The boxed neighborhood supplies an aggregate state to the focal agent's rule; heterogeneous agents may use different rule tables on the same graph.}
  \label{fig:wolfram-grid}
\end{figure}

\section{Synchronous Dynamics and Boundary Experts}

With synchronous updating, every agent reads the same frame $\mathbf s(t)$ and all next states are written to a new frame $\mathbf s(t+1)$. Updating the grid in place makes later agents respond to already-updated neighbors and therefore implements an order-dependent asynchronous model. Neither scheme is universally correct, but they are different models and should not be mixed accidentally.

The course notebook uses boundary agents as ``experts.'' A moving typical-price signal can assign their action using a neutral band:
\begin{equation}
  \label{eq:technical-boundary}
  s_{\mathrm{edge}}(t+1)=
  \begin{cases}
    0, & X_t<\overline X_t-\beta\widehat\sigma_t,\\
    1, & |X_t-\overline X_t|\leq\beta\widehat\sigma_t,\\
    2, & X_t>\overline X_t+\beta\widehat\sigma_t,
  \end{cases}
\end{equation}
where $X_t$ is a declared price statistic, $\overline X_t$ a trailing mean, $\widehat\sigma_t$ a trailing scale, and $\beta\geq0$ controls the hold region. The labels ``buy'' below trend and ``sell'' above trend encode a contrarian rule. Reversing them produces momentum. The economic interpretation must match the inequality.

If all boundary agents share \eqnref{technical-boundary}, their simultaneous signal is an intentionally strong common shock. Interior agents may imitate a majority rule, follow a minority rule, or use heterogeneous $\phi_i$. Heterogeneity can damp, fragment, or amplify propagation depending on the rules and topology; it is not by itself a theorem of instability.

\notebooklink
  {L15a: Wolfram Network Simulation}
  {https://github.com/varnerlab/CHEME-5660-CourseRepository-Fall-2026/blob/main/lectures/week-15/L15a/CHEME-5660-L15a-Example-Wolfram-NetworkSimulation-Fall-2025.ipynb}
  {Construct moving-price boundary signals, encode totalistic policies, and simulate signal propagation through connected agents.}

\notebooklink
  {L15a: Heterogeneous-Agent Crash Simulation}
  {https://github.com/varnerlab/CHEME-5660-CourseRepository-Fall-2026/blob/main/lectures/week-15/L15a/CHEME-5660-L15a-Example-Wolfram-NetworkSimulation-Crash-Fall-2025.ipynb}
  {Combine majority and minority policies and inspect whether sell-state cascades emerge under a selected ruleset.}

\section{A State Movie Is Not Yet a Market}

Colored cell dynamics visualize action propagation, but a market model also needs order size, inventory, clearing, and price formation. One simple bridge maps actions to signed demand $g(0)=+1$, $g(1)=0$, and $g(2)=-1$ and defines normalized imbalance
\begin{equation}
  \label{eq:imbalance}
  Q_t:=\frac{1}{N}\sum_{i=1}^{N}v_i(t)g(s_i(t)),
\end{equation}
where $v_i(t)\geq0$ is agent $i$'s order scale. A stylized impact equation is
\begin{equation}
  \label{eq:price-impact}
  r_{t+1}=\lambda_t Q_t+\sigma_t\varepsilon_{t+1},
  \qquad
  P_{t+1}=P_t\exp(r_{t+1}),
\end{equation}
where $\lambda_t>0$ is inverse liquidity, $\sigma_t\geq0$ sets exogenous noise scale, and $\varepsilon_{t+1}$ has zero mean and unit variance. Fixed $v_i$ permits unlimited repeated buying or selling; a credible extension enforces cash, inventory, leverage, and market-clearing constraints.

Endogenous feedback occurs only if the resulting $P_{t+1}$ changes subsequent signals, constraints, or rules. Precomputing boundary states from an observed QQQ path is a driven network experiment: prices influence agents, but agents do not influence that historical path. It can study propagation, not demonstrate that those agents caused the observed market moves.

Evaluate an agent-based market at three levels:
\begin{itemize}[leftmargin=1.3em,itemsep=2pt,topsep=2pt]
  \item \textbf{Verification:} Does the rule index reproduce its table, are boundaries and update order correct, and are results reproducible from seed and parameters?
  \item \textbf{Validation:} Do simulated returns match several out-of-sample targets---tails, volatility clustering, drawdowns, volume, and recovery duration---rather than one attractive crash path?
  \item \textbf{Identification:} Are conclusions stable across plausible rules, networks, initial states, liquidity maps, and parameter perturbations?
\end{itemize}
Matching stylized facts is evidence of generative adequacy for those facts, not proof that the encoded mechanism caused a historical crash.

\keytakeaways
  {In this lecture, we connected crash-amplification mechanisms to a transparent agent-based market built from local three-state Wolfram rules.}
  {Crashes combine paths and feedback}
  {Returns and drawdowns measure different features, while leverage, liquidity, crowding, contagion, and reflexive constraints can amplify an initial shock.}
  {Update conventions define the model}
  {Totalistic rules collapse neighbor configurations to a sum; rule-digit order, boundaries, topology, and synchronous versus asynchronous updates change the dynamics.}
  {Emergence is not causal proof}
  {Agent states need an inventory and price-formation layer, and reproducing a crash-like pattern must be followed by robustness and out-of-sample validation.}
  {Next, integrate the course models into an auditable 2026 Autotrader launch pipeline.}

\begin{readinglist}
  \item Alan Kirman, \href{https://doi.org/10.2307/2233844}{``Ants, Rationality, and Recruitment,''} \emph{Quarterly Journal of Economics} 108(1), 137--156 (1993).
  \item Thomas Lux and Michele Marchesi, \href{https://doi.org/10.1038/17290}{``Scaling and Criticality in a Stochastic Multi-Agent Model of a Financial Market,''} \emph{Nature} 397, 498--500 (1999).
  \item SEC and CFTC, \href{https://www.sec.gov/files/marketevents-report.pdf}{\emph{Findings Regarding the Market Events of May 6, 2010}} (2010).
\end{readinglist}

{\sffamily\footnotesize\color{vnmuted}\textbf{Educational use and risk notice.} These notes are for instruction only. A simulated crash is not a forecast, and an agent-based model is not investment advice.}

\end{document}
